\section{Methodology}

To investigate the effectiveness of the proposed Jupyter dashboard extension in supporting student reflection and reducing programming errors, a mixed-methods research design will be employed. Our study will mainly focus on answering the following research questions:
\begin{enumerate}
    \item[RQ1:] Does the inclusion of a learner-facing dashboard in a Jupyter-based programming environment help students better understand and reflect on their own learning process?
    \item[RQ2:] How do students perceive the usefulness and usability of the dashboard for identifying and correcting mistakes?
    \item[RQ3:] Which usage indicators can serve as signals for frustration, reflection or effective learning strategies?
\end{enumerate}

In order to address these questions, our study will mainly consist of two phases: a first phase focused on starting using the dashboard in real classroom settings, and a second phase involving a more advanced notebook whose aim is the reflection of the students' learning process. In order to evaluate the impact of the dashboard, we will divide our participants into two main groups:
\begin{itemize}
    \item \underline{Control Group}: Participants in this group will use Jupyter Notebooks without access to the dashboard features. They will complete programming exercises as usual, without any additional feedback or visualizations.
    \item \underline{Experimental Group}: Participants in this group will have access to the Jupyter dashboard extension while completing the same programming exercises. They will be encouraged to use the dashboard to monitor their activity, reflect on their errors, and adjust their strategies accordingly.
\end{itemize}
% // TODO: Maybe explain more from where the participants will be recruited.

It is important to note that both groups will work on the same programming tasks and exercises, and that the group divission will be randomized to minimize selection bias, ensuring that any differences in outcomes can be attributed to the presence or absence of the dashboard. Our study design is divided in a defined and structured number of steps, which are described in the following sections.

\subsection{Initial Survey}

Before starting the learning units, all participants will complete a pre-survey to collect prior programming experience and baseline attitudes towards programming and self-assessment. This survey will help us control for potential confounding variables and understand the initial context of our participants. This information will allow us to analyze whether the dashboard's effectiveness varies based on students' initial experience levels or attitudes.
Regarding instructors, we will also conduct a brief survey to gather how they usually grade student's notebooks, what difficulties they face when assessing student work, and what feedback (if any) they provide to students about their coding process.

This first survey will set the stage for our subsequent interventions and analyses. It will provide valuable context for interpreting the results of the learning units and will also help us tailor the dashboard features to better meet the needs of our participants.

It should be highlighted that no demographic data will be collected in order to preserve the anonymity of the participants, so no ethic approval will be needed for this study.

\subsection{First Learning Unit}

The first learning unit will introduce participants to basic Python programming concepts, such as variables, data types, and simple expressions. Participants will complete a series of programming exercises within Jupyter Notebooks. As mentioned earlier, the control group will not have access to the dashboard, while the experimental group will be able to use it to monitor their activity. Depending of the instructors' preferences, the experimental group students will also be given some feedback according to the information shown in the professor's dashboard. This feedback may include hints about common errors, suggestions for improvement, or prompts to reflect on their coding strategies.

The content of both dashboards will be designed depending on the results of the initial survey, but it may include metrics such as the total number of executed cells, error rates or even a timeline of activity.

\subsection{Second Learning Unit}

The second learning unit will build upon the concepts introduced in the first unit, introducing slightly more advanced topics such as loops and conditionals. Participants will again complete programming exercises within Jupyter Notebooks. The experimental group will continue to have access to the dashboard, allowing them to review their previous performance and errors from the first unit. This will enable us to examine whether learners who could review their behavior and performance in the first unit make fewer errors and perform better in the second unit.

\subsection{Post-Survey}

After completing both learning units, all participants will complete a post-survey to assess their experiences with the dashboard (for the experimental group) and their overall perceptions of the learning process. The survey will include Likert-scale items and open-ended questions to gather both quantitative and qualitative data, and will focus on aspects such as the perceived usefulness of the dashboard, its impact on self-reflection, and any changes in learning strategies.

After collecting and analyzing the data from both learning units and the surveys, we will be able to draw conclusions about the effectiveness of the Jupyter dashboard extension in supporting student reflection and reducing programming errors. The findings will inform future iterations of the dashboard and contribute to our understanding of how technology can enhance programming education.



% // I have not explained the contents of the dashboards in detail, as they may vary depending on the results of the initial survey. Once we have that information, we can decide which metrics and visualizations will be most relevant for the students and instructors.